%=====================================================================
% Sección de Dorantes Mejía Emmanuel: bitácora clase 21 (drivers)
%=====================================================================
\subsection{Clase 21: Driver de cada sensor y actuador}
\textbf{Autor:} Dorantes Mejía Emmanuel

\textbf{Pregunta:} Para cada sensor o actuador de su proyecto,
indiquen qué archivo de dispositivo o biblioteca usa y qué
driver del kernel hay detrás (con \texttt{ls -l /dev},
\texttt{lsmod} o \texttt{dmesg} como evidencia).
Si alguna parte necesitara un driver propio, expliquen
por qué no alcanza el genérico.

\medskip
\textbf{Respuesta:}

El proyecto utiliza una Raspberry Pi Pico para adquirir las
señales de entrada y controlar los actuadores, mientras que
la Raspberry Pi 4 ejecuta Linux y gestiona la comunicación
USB y la visualización del videojuego.

En la Pico, los periféricos se controlan mediante el SDK
y código propio, sin utilizar drivers del kernel Linux.
En la Raspberry Pi 4 se emplean los controladores
correspondientes a los dispositivos conectados.

\begin{table}[H]
    \centering
    \caption{Drivers de sensores y actuadores}
    \label{tab:drivers}
    \small
    \begin{tabular}{|p{1.8cm}|p{2.6cm}|p{2.6cm}|}
        \hline
        \textbf{Elemento} & \textbf{Driver} &
        \textbf{Evidencia} \\
        \hline
        Botones &
        PIO + DMA (propio) &
        \texttt{botones.c} y
        \texttt{read\_74hc165.pio} \\
        \hline
        Potenciómetro &
        \texttt{hardware/adc} y
        \texttt{hardware/dma} &
        \texttt{read\_adc.c} y
        \texttt{read\_adc.h} \\
        \hline
        Motor, buzzer &
        \texttt{hardware/pwm} &
        \texttt{actuadores.c} \\
        \hline
        LEDs &
        \texttt{hardware/gpio} &
        \texttt{actuadores.c} \\
        \hline
        Control (Pi) &
        \texttt{usbhid}
        $\rightarrow$
        \texttt{/dev/input/eventX} &
        \texttt{/proc/bus/input/devices}
        (pendiente de verificar) \\
        \hline
        Serial (Pi) &
        \texttt{cdc\_acm}
        $\rightarrow$
        \texttt{/dev/ttyACM0} &
        \texttt{dmesg}
        (pendiente de verificar) \\
        \hline
        Pantalla (Pi) &
        \texttt{panel-mipi-dbi} &
        \texttt{lsmod}
        (pendiente de verificar) \\
        \hline
    \end{tabular}
\end{table}

\medskip
\textbf{Aportación individual:}

Se desarrollaron los archivos \texttt{read\_adc.c} y
\texttt{read\_adc.h} para implementar la adquisición
continua de señales analógicas mediante el ADC de
12 bits de la Raspberry Pi Pico.

La implementación configura los canales ADC0, ADC1,
ADC2 y ADC3 mediante muestreo round-robin, con una
frecuencia configurada de 4000 muestras por segundo
por canal.

Se utilizan dos canales DMA encadenados para transferir
automáticamente las muestras hacia un búfer circular
de cuatro elementos de 16 bits, reduciendo la
intervención de la CPU.

El archivo \texttt{read\_adc.c} contiene la inicialización
del ADC, la configuración del DMA y la función
\texttt{read\_adc\_is\_alive()}, que permite supervisar
el avance de las transferencias.

Por su parte, \texttt{read\_adc.h} define las constantes
de configuración y declara las funciones utilizadas
por el módulo.

Esta implementación utiliza los periféricos ADC y DMA
del microcontrolador mediante las bibliotecas del
Pico SDK, sin requerir un driver del kernel Linux.

Para la lectura de botones se utiliza código propio
basado en PIO y DMA, debido a que se requiere controlar
la señal de carga paralela y el reloj de los registros
74HC165. Los demás periféricos emplean las bibliotecas
y controladores indicados en la tabla.

\begin{figure}[H]
    \centering
    \imgph{Capturas pendientes de
    \texttt{lsusb},
    \texttt{dmesg | tail},
    \texttt{cat /proc/bus/input/devices}
    y \texttt{lsmod | grep mipi}}{4cm}
    \caption{Verificación de los controladores
    de dispositivos en la Raspberry Pi 4.}
    \label{fig:drivers}
\end{figure}
